% Pursuing Stacks (cotes 134-2 to 134-8): archive cote 134-3, site page 75 (PDF page 76), typescript page 75.
% typescript. Draft compared: ps2.tex:3875-3949.
% AI-transcribed by Claude Opus 5.5 (claude-opus-5-5, adaptive thinking, effort high), Message Batches API, batch ps1 (msgbatch_01Rv5hCaNTtUt4JHx1yqoHPL), 2026-11-12; not yet checked by eye.
\scanpage{134-3/75}
\mspage{75}

43.   Now let's come back for a little while again to the so called ``standard
test categories'', and check how nicely the ``story retold'' applies to them.

Not speaking about multicomplexes, there are essentially two variants
for ``categories of simplices'' as test categories. The smaller, more comm-%
monly used one, is the category $\hand{\Delta}$ of ``ordered simplices'', most convenient%
ly described as the full subcategory of (Cat) defined by the family of
simplices $\hand{\Delta}_n$ ($n \hand{\in} \underline{N}$). Here the most natural choice for i: $\hand{\Delta} \hand{\longrightarrow}$ (Cat) is
of course the inclusion functor. \deleted{The criterion for what might be called
``test-functors'' applies trivially here (cor. 3 of ``key theorem'', p.66),
\inserted[typed]{provided we check that}
thus \quad is indeed a test category and i a test-functor, \xout{\unsure{moreover}} \quad$_1$ to-%
gether with the usual \quad$_o$, \quad$_1$ is a homotopy interval in \quad. \xout{\unsure{Also, the stan-}}
\xout{\unsure{dard criterion}}} As i is fully faithful and $\hand{\Delta}_1$ is in the image, it follows
that $i^{\hand{*}}(\hand{\Delta}_1) = \hand{\Delta}_1$ , and we have only to check (for A to be a test-category
with ``test-functor'' i) that $\hand{\Delta}_1$ is \xout{an} aspheric over $e=\hand{\Delta}_o$ , namely that
all products $\hand{\Delta}_1 \hand{\times} \hand{\Delta}_n$ are aspheric - which we did. The extra condition
of the``total asphericity criterion'' (proposition on \over{\unsure{p. 67}}{\S. 39}), namely existen%
\hand{-}ce of a homotopy \inserted{in $\Delta^{\wedge}$} from the \xout{constant} identity map to a constant map\hand{,} for any
$\hand{\Delta}_n$ , is indeed satisfied: \xout{\illegible{1}}it \xout{\unsure{amounts to the same}} \inserted{is enough} to define such homotopy
in (Cat), which is trivial using the final element of $\hand{\Delta}_n$ . Thus $\hand{\Delta}$ is
in fact a \uline{strict} test category.

As for $\hand{\widetilde{\Delta}}$ , the most elegant choice theoretically is to take the cate-%
gory of all non-empty finite sets, but this leads to set-theoretic diffi-%
culties, as this category is not small - thus we take again the standard
non-ordered simplices
\[ \hand{\widetilde{\Delta}}_n = \underline{N} \hand{\cap} \hand{[}0,n\hand{]} \qquad (n \hand{\in} \underline{N}) \quad , \]
so as to get a ``reduced'' category with a countable set of objects. This
time, as $\hand{\widetilde{\Delta}}$ is stable under finite products, and contains the ``interval''
$(\hand{\widetilde{\Delta}}_1, \hand{\delta}_o, \hand{\delta}_1)$ (which is necessarily then a homotopy interval, as all ele-%
ments of $\hand{\widetilde{\Delta}}$ are aspheric over $e = \hand{\widetilde{\Delta}}_o$), the fact that $\hand{\Delta}$ is a strict test
category is trivial. As for a test functor, the neatest choice is the one
we said before, namely associating to every finite set the ordered set of
all non-empty subsets. We thus get
\[ \hand{\widetilde{\ }}i : \hand{\widetilde{\Delta}^{\unsure{\wedge}}} \hand{\longrightarrow} (Cat) \qquad \text{(factoring in fact through (Ord),} \]
\[ \text{as does } i: \hand{\Delta} \hand{\longrightarrow} \text{(Cat) above )} \]
To prove it is indeed a test functor, the corollary 3 to the ``key theorem''
(\over{\unsure{p.66}}{\S.38}) applies, taking of course I = $\hand{\widetilde{\Delta}}_1$ , hence
\[ \hand{\widetilde{\ }}i_{\hand{!}} (I) = i(\hand{\widetilde{\Delta}}_1 ) = \hand{\diagram{\begin{tikzcd}[row sep=small, column sep=small] \{1\} \ar[dr] & \\ & \{0,1\} \\ \{0\} \ar[ur] & \end{tikzcd}}} \quad . \]
